\chapter{Boolean Analysis — Restriction Monotonicity}
%%% generated by the blueprint pipeline from TCSlib.BooleanAnalysis.LMN.RestrictionMonotonicity
\section{Overview}

This module shows that decision-tree depth can only decrease when a function is further
restricted. It builds the operation of restricting a decision tree by a partial
assignment, proves the resulting tree is shallower and computes the restricted function,
and derives the key Bernoulli monotonicity $\Pr_{p\cdot q}[\mathrm{dtDepth} > t] \le
\Pr_{p}[\mathrm{dtDepth} > t]$.

%%%%%%%%%%%%%%%%%%%%%%%%%%%%%%%%%%%%%%%%%%%%%%%%%%%%%%%%%%%%%%%%%%%%%%%%%%%%%%%
\section{Declarations}

\begin{definition}[Restriction of a decision tree]
\lean{LMN.dtRestrict}
\label{LMN.dtRestrict}
\leanok
\uses{DecisionTree, SwitchingLemma2.Restriction}
Given a decision tree $T$ and a partial assignment $\rho$, the restricted tree
$\mathrm{dtRestrict}\,T\,\rho$ is obtained by collapsing every branch whose query variable
is fixed by $\rho$ into the corresponding subtree, and keeping unrestricted branches intact.
\end{definition}

\begin{theorem}[Restriction does not increase decision-tree depth]
\lean{LMN.dtRestrict_depth_le}
\label{LMN.dtRestrict_depth_le}
\leanok
\uses{DecisionTree, DecisionTree.depth, LMN.dtRestrict, SwitchingLemma2.Restriction}
\difficulty{4}
Let $T$ be a decision tree on $n$ Boolean variables, and let $\rho$ be a restriction on
those variables, that is, an assignment giving each variable either a fixed Boolean
value or leaving it free. Let $T|_\rho$ denote the tree obtained by collapsing every
branch whose query variable is fixed by $\rho$ into the corresponding subtree while
keeping the unrestricted branches intact, and write $\mathrm{depth}$ for the maximum
length of a path from the root to a leaf. Then
\[
\mathrm{depth}(T|_\rho) \le \mathrm{depth}(T).
\]
\end{theorem}

\begin{theorem}[Restricting a decision tree commutes with evaluation]
\lean{LMN.dtRestrict_eval}
\label{LMN.dtRestrict_eval}
\leanok
\uses{DecisionTree, DecisionTree.eval, LMN.dtRestrict, SwitchingLemma2.Restriction, SwitchingLemma2.Restriction.extend}
\difficulty{4}
Let $T$ be a decision tree on $n$ Boolean variables, let $\rho$ be a restriction on
those variables, and let $x \in \{\mathrm{true},\mathrm{false}\}^n$ be an input. Then
evaluating the restricted tree at $x$ agrees with evaluating $T$ at the assignment
obtained by extending $\rho$ with $x$; that is, the restricted tree returns the same bit
at $x$ as $T$ does at the point that takes each variable's fixed value under $\rho$
where $\rho$ sets it and the value $x_i$ on the free variables.
\end{theorem}

\begin{theorem}[Decision-tree depth is monotone under restriction]
\lean{LMN.dtDepth_restrictFn_le'}
\label{LMN.dtDepth_restrictFn_le'}
\leanok
\uses{DecisionTree, DecisionTree.depth, DecisionTree.eval, LMN.dtRestrict, LMN.dtRestrict_depth_le, LMN.dtRestrict_eval, SwitchingLemma2.Restriction, SwitchingLemma2.restrictFn, buildFullDTree, buildFullDTree_depth, buildFullDTree_eval, dtDepth}
\difficulty{5}
Let $f$ be a Boolean function of $n$ Boolean variables, and let $\rho$ be a restriction
that fixes some of these variables to Boolean values and leaves the remaining ones free.
Write $f_\rho$ for the restricted function, whose value on an input $x$ is $f$ evaluated
at the assignment that agrees with $\rho$ on the fixed coordinates and with $x$ on the
free coordinates. Writing $D(g)$ for the decision-tree depth of a Boolean function $g$ —
the least depth of a decision tree computing $g$ — one has $D(f_\rho) \le D(f)$.
\end{theorem}

\begin{theorem}[Composing restrictions does not increase decision-tree depth]
\lean{LMN.dtDepth_composeRestr_le}
\label{LMN.dtDepth_composeRestr_le}
\leanok
\uses{LMN.composeRestr, LMN.dtDepth_restrictFn_le', SwitchingLemma2.Restriction, SwitchingLemma2.Restriction.extend, SwitchingLemma2.restrictFn, dtDepth}
\difficulty{4}
Let $f$ be a Boolean function of $n$ variables, and let $\rho_1$ and $\rho_2$ be two
restrictions on those variables, each pinning some coordinates to a fixed Boolean value
and leaving the rest free. Let $\rho_1 \circ \rho_2$ be the composed restriction that
fixes each coordinate $i$ to the value assigned by $\rho_1$ when $\rho_1$ pins $i$, and
otherwise to the value assigned by $\rho_2$. For a restriction $\rho$, write $f|_{\rho}$
for the function obtained from $f$ by substituting the fixed values of $\rho$ and
treating the coordinates left free by $\rho$ as its inputs, and let
$\operatorname{depth}_{\mathrm{DT}}(g)$ denote the least $d$ for which some decision
tree of depth at most $d$ computes $g$. Then further restricting along $\rho_2$ can only
lower this depth:
\[
\operatorname{depth}_{\mathrm{DT}}\bigl(f|_{\rho_1 \circ \rho_2}\bigr) \le
\operatorname{depth}_{\mathrm{DT}}\bigl(f|_{\rho_1}\bigr).
\]
\end{theorem}

\begin{theorem}[Monotonicity of the decision-tree depth tail under Bernoulli restrictions]
\lean{LMN.bernoulliRestrProb_dtDepth_compose_le}
\label{LMN.bernoulliRestrProb_dtDepth_compose_le}
\leanok
\uses{LMN.composeRestr, LMN.dtDepth_composeRestr_le, LMN.restriction_compose_eq, SwitchingLemma2.bernoulliRestrProb, SwitchingLemma2.bernoulliRestrProb_le_one', SwitchingLemma2.bernoulliRestrWeight, SwitchingLemma2.bernoulliRestrWeight_nonneg', SwitchingLemma2.restrictFn, dtDepth}
\difficulty{5}
Let $f : (\mathrm{Fin}\,n \to \mathrm{Bool}) \to \mathrm{Bool}$ be a Boolean function,
let $t \in \bbn$ be a threshold, and let $p_1, p_2 \in \bbr$ satisfy $0 < p_1 \le 1$ and
$0 < p_2 \le 1$. For a parameter $p$, let $\Pr_p[\,\cdot\,]$ denote the probability of
an event under a Bernoulli($p$) random restriction $\rho$, in which each variable is
left free with probability $p$ and is otherwise fixed to a constant Boolean value, each
of the two values with probability $(1-p)/2$. For a restriction $\rho$, write $f_\rho$
for the function of the free coordinates obtained from $f$ by fixing each set coordinate
$i$ to the value $\rho(i)$, and let $D(f_\rho)$ be its decision-tree depth, that is, the
least $d$ for which some decision tree of depth at most $d$ computes $f_\rho$. Then
\[
  \Pr_{p_1 p_2}\bigl[D(f_\rho) > t\bigr] \le
  \Pr_{p_1}\bigl[D(f_\rho) > t\bigr].
\]
\end{theorem}
